\mxnewpage\section*{Correction EV3-T2}\label{EV3-T2-CORR}
\textbf{QCM :} 1 D, 2 D, 3 B, 4 D, 5 C, 6 C, 7 A, 8 A, 9 B, 10 A, 11 B, 12 A. Chaque réponse correcte : \num{0,5} point ; total 6.
\paragraph{Vérifications des automatismes.} Priorités ; somme de fractions ; soustraction des exposants ; coefficient multiplicateur \num{1,2} ; division par \num{0,8} ; produit \num{1,25} fois \num{0,8} ; résolution affine ; accroissement vertical divisé par l’horizontal ; moyenne ; angle divisé par 360 ; complément à 1 ; substitution entre parenthèses.
\subsection*{Exercice 1 -- Simulation et moyenne}
\begin{enumerate}
\item $E(X)=0+\dfrac{1}{2}+\dfrac{3}{2}=2$. \textit{Barème : 1 point. Résultat et justification élémentaire attendus ; demi-point possible pour démarche correcte inachevée.}
\item $E(X^2)=1+9=10$, variance 6, écart type $\sqrt6$. \textit{Barème : 2 points. Méthode ou justification : 1 ; calcul et conclusion : 1.}
\item $\intff{2-\dfrac{2\sqrt 6}{5}}{2+\dfrac{2\sqrt 6}{5}}$. Une fréquence sur un nombre fini de simulations fluctue. \textit{Barème : 2 points. Méthode ou justification : 1 ; calcul et conclusion : 1.}
\end{enumerate}
\subsection*{Exercice 2 -- Comparer deux modèles}
\begin{enumerate}
\item $h\prime=2(e^{2x}-1)$. \textit{Barème : 1 point. Résultat et justification élémentaire attendus ; demi-point possible pour démarche correcte inachevée.}
\item Décroît jusqu’à 0 puis croît ; minimum 0. Donc $f\geq g$ avec égalité seulement en 0. \textit{Barème : 2 points. Méthode ou justification : 1 ; calcul et conclusion : 1.}
\item Raison $e^2$ ; somme $1+e^2+e^4+e^6=\dfrac{1-e^8}{1-e^2}$. \textit{Barème : 2 points. Méthode ou justification : 1 ; calcul et conclusion : 1.}
\end{enumerate}
\subsection*{Exercice 3 -- Un cercle déplacé}
\begin{enumerate}
\item Centre (1;-1), rayon 2. \textit{Barème : 1 point. Résultat et justification élémentaire attendus ; demi-point possible pour démarche correcte inachevée.}
\item A(-1;-1), B(3;-1). \textit{Barème : 1 point. Résultat et justification élémentaire attendus ; demi-point possible pour démarche correcte inachevée.}
\item AB est un diamètre, donc produit 0, extrémités comprises ; identité $MI^2-4=0$. \textit{Barème : 2 points. Méthode ou justification : 1 ; calcul et conclusion : 1.}
\end{enumerate}
\paragraph{Étalonnage.} Durée cible : 20–25 min pour les automatismes, 80–85 min pour les exercices, 10–15 min de relecture. Ces durées sont des estimations éditoriales à mesurer en classe. Accepter toute méthode correcte ; ne pas sanctionner deux fois une erreur numérique propagée avec cohérence. Le barème proposé ne vaut pas décision d’harmonisation de l’équipe.
